\section{Literature Review}

\subsection{Analytical Orbit Propagation}
The foundational work on analytical orbit propagation was established by Brouwer (1959) and later refined by Lane and Cranford (1969), leading to the development of the SGP4 algorithm by Hoots and Roehrich (1980). SGP4 utilizes a simplified analytical model to account for Earth's oblateness ($J_2, J_3, J_4$ harmonics), third-body perturbations (luni-solar effects), and a static atmospheric density model. While computationally superior, SGP4's inability to ingest real-time space weather parameters fundamentally limits its covariance realism, especially for prediction horizons exceeding 12 hours in LEO.

\subsection{Machine Learning in Astrodynamics}
The integration of Machine Learning (ML) to correct orbital propagation errors has gained massive traction. Early works utilized Support Vector Machines (SVMs) and Random Forests to predict the degradation of TLEs. More recently, Recurrent Neural Networks (RNNs) and LSTMs have been utilized to capture the time-series degradation of the SGP4 residual error. However, as noted by several researchers, these models suffer from ''black-box'' unreliability. They often fail to extrapolate outside their training distributions and suffer from initial-state violations, severely limiting their integration into automated safety-of-flight systems.

\subsection{Physics-Informed Neural Networks (PINNs)}
Physics-Informed Neural Networks (PINNs), introduced by Raissi et al. (2019), have revolutionized computational science by incorporating physical governing equations (e.g., Navier-Stokes, orbital mechanics) directly into the loss function of a neural network as soft constraints. In astrodynamics, PINNs have been explored for gravity field modeling and optimal control. However, using soft constraints (loss penalties) to enforce the initial state vector at $t=0$ often fails in highly chaotic systems like LEO drag, as the optimizer trades off initial-state accuracy to minimize long-term prediction loss. Our work bypasses this limitation by proposing a ''hard constraint'' gating layer, guaranteeing absolute physical compliance at the epoch while retaining the flexibility of a deep neural architecture for future time steps.

\subsection{Space Weather and Drag Prediction}
The correlation between space weather, thermospheric density, and satellite drag is well-documented. Density models like NRLMSISE-00 and JB2008 incorporate F10.7, $S_{10.7}$, and $K_p$ indices to estimate density. High Precision Orbit Propagators (HPOP) utilize these density models for numerical integration. However, numerical integration is too slow for all-on-all conjunction assessments (involving millions of pairs). Our G-PINN acts as a surrogate model, effectively bridging the speed of SGP4 with the environmental awareness of an HPOP by injecting Celestrak F10.7 observations directly into the latent space of the neural network.
